\documentclass[aspectratio=169,11pt]{beamer}

\usepackage{iftex}
\ifXeTeX
  \usepackage{fontspec}
  \usepackage{xeCJK}
  \setmainfont{Times New Roman}
  \IfFontExistsTF{Microsoft JhengHei}{
    \setCJKmainfont{Microsoft JhengHei}
    \setCJKsansfont{Microsoft JhengHei}
    \setCJKmonofont{Microsoft JhengHei}
  }{
    \setCJKmainfont{SimSun}
    \setCJKsansfont{SimSun}
    \setCJKmonofont{SimSun}
  }
  \IfFontExistsTF{Consolas}{\setmonofont{Consolas}}{}
\fi
\usepackage{amsmath}
\usepackage{array}
\usepackage{booktabs}
\usepackage{graphicx}
\usepackage{siunitx}
\usepackage{xcolor}
\usepackage{xurl}

\usetheme{Madrid}
\usecolortheme{default}
\setbeamertemplate{navigation symbols}{}
\setbeamertemplate{footline}[frame number]
\graphicspath{{images/baseline_comparison_20260509/}{images/lono_20260509/}}
\urlstyle{same}

\definecolor{aaltoBlue}{RGB}{0,83,155}
\definecolor{goodGreen}{RGB}{35,120,75}
\definecolor{warnRed}{RGB}{170,55,55}
\definecolor{softGray}{RGB}{245,246,248}

\newcommand{\good}[1]{\textcolor{goodGreen}{#1}}
\newcommand{\warn}[1]{\textcolor{warnRed}{#1}}
\newcommand{\blue}[1]{\textcolor{aaltoBlue}{#1}}
\newcommand{\sourcepath}[1]{{\scriptsize\url{#1}}}
\newcommand{\sourcenote}[1]{\vfill{\tiny\textcolor{gray!70!black}{Source: \url{#1}}}}
\newcolumntype{L}[1]{>{\raggedright\arraybackslash}p{#1}}
\newcolumntype{R}[1]{>{\raggedleft\arraybackslash}p{#1}}

\title[Baseline Comparison]{Baseline 对比与 \texorpdfstring{\(\Delta P^{0.5}\)}{Delta P exponent 0.5} Stage-3 MLP}
\subtitle{Hiroyasu--Arai、Naber--Siebers、Sparse GP 与五 seed 外部 clean eval}
\author{Linan Jiang}
\institute{Aalto University}
\date{9 May 2026}

\begin{document}

\begin{frame}
  \titlepage
  \vspace{-0.8em}
  \centering
  \tiny Frozen slide data snapshot:
  \url{Thesis/generated/baseline_comparison_20260509/}
\end{frame}

\begin{frame}{这套 slides 的证据链}
\small
\begin{tabular}{L{0.24\linewidth}L{0.33\linewidth}L{0.34\linewidth}}
\toprule
问题 & 方法 & 读法 \\
\midrule
物理公式能否解释数据？ &
Hiroyasu--Arai 与 Naber--Siebers &
确定性物理相关式；给出强但刚性的 reference。 \\
普通 ML baseline 是否已足够？ &
Sparse heteroscedastic GP &
非参数、带不确定度；检验 MLP 是否只是吃到小数据优势。 \\
神经网络是否稳定？ &
\(\Delta P^{0.5}\) Stage-3，5 seed &
看 seed-to-seed 波动，而不是单次 run headline。 \\
是否泛化到未见喷嘴？ &
Leave-one-nozzle-out &
整组喷嘴从训练中移除，测试真正的 OOD 插值/外推压力。 \\
不确定度是否可信？ &
Coverage、reliability、CRPS、PIT &
区分“误差小”和“误差条有统计意义”。 \\
\bottomrule
\end{tabular}

\vspace{0.5em}
\footnotesize
本 deck 的重点不是只报一个 RMSE，而是把 accuracy、tail error、calibration、
seed stability 与 OOD behaviour 放在同一协议框架下解释。
\end{frame}

\begin{frame}{指标含义：accuracy、tail 与 calibration}
\small
\begin{columns}[T,onlytextwidth]
\begin{column}{0.48\linewidth}
\begin{block}{点预测误差}
\footnotesize
\[
\mathrm{RMSE}=\sqrt{\frac{1}{n}\sum_i(\hat y_i-y_i)^2},
\quad
\mathrm{MAE}=\frac{1}{n}\sum_i|\hat y_i-y_i|.
\]
RMSE 对大误差更敏感；MAE 更接近日常意义上的平均毫米误差。
\[
\mathrm{Bias}=\frac{1}{n}\sum_i(\hat y_i-y_i)
\]
用于判断系统性高估或低估。
\end{block}
\end{column}

\begin{column}{0.48\linewidth}
\begin{block}{尾部误差与不确定度}
\footnotesize
\[
P95=\mathrm{quantile}_{0.95}(|\hat y-y|)
\]
表示 95\% 点的误差不超过该值，直接反映 worst-case tail。
\[
\mathrm{Coverage}_{k\sigma}
=\Pr(|y-\mu|\leq k\sigma).
\]
若 Gaussian 预测校准良好，1\(\sigma\)/2\(\sigma\) 应接近
0.683/0.954。高于 nominal 是 over-coverage，通常更保守但不够 sharp。
\end{block}
\end{column}
\end{columns}
\end{frame}

\begin{frame}{协议含义：in-pipeline vs external clean}
\small
\begin{block}{In-pipeline 5-seed bootstrap}
\footnotesize
每个 seed 完整重跑 Stage-1/2/3 pipeline，并在该 seed 的 post-evaluation
test subset 上计算 winner metrics。Bootstrap 的对象是 seed-level metric：
对五个 seed 有放回重采样，得到均值的经验分布与 95\% interval。
它回答的是：\textbf{训练随机性会不会改变结论？}
\end{block}

\begin{block}{External clean protocol}
\footnotesize
训练好的模型再统一评估到同一个 full-clean diagnostic set：
\num{795561} points，\num{33845} trajectories。它回答的是：
\textbf{所有方法在同一大规模外部 clean 集合上谁更准？}
\end{block}

\begin{block}{严谨表述}
\footnotesize
这两个协议不能直接混合成一个数字。本文把 in-pipeline 结果用于 seed stability，
把 external clean 结果用于 headline model comparison。
\end{block}
\end{frame}

\begin{frame}{一句话结论}
\small
\begin{block}{同一外部 clean eval 协议}
\(\Delta P^{0.5}\) 特征工程后的 Stage-3 MLP 五 seed 均值，在
\num{795561} 个点、\num{33845} 条轨迹上达到
\good{RMSE \SI{6.211}{mm} \(\pm\) \SI{0.283}{mm}}。
\end{block}

\vspace{0.4em}
\begin{columns}[T,onlytextwidth]
\begin{column}{0.48\linewidth}
\begin{block}{相对物理 baseline}
\begin{itemize}
  \item vs Hiroyasu--Arai：RMSE \good{-37.7\%}，MAE \good{-40.3\%}。
  \item vs Naber--Siebers：RMSE \good{-29.7\%}，MAE \good{-32.3\%}。
  \item vs sparse GP：RMSE \good{-10.0\%}，MAE \good{-11.9\%}。
\end{itemize}
\end{block}
\end{column}

\begin{column}{0.48\linewidth}
\begin{block}{相对旧 MLP}
\begin{itemize}
  \item 旧 \(\Delta P^{0.25}\) seed 42：RMSE \SI{6.696}{mm}。
  \item 新 \(\Delta P^{0.5}\) 五 seed 均值：RMSE \SI{6.211}{mm}。
  \item Sparse GP mean：RMSE \SI{6.898}{mm}，MAE \SI{5.206}{mm}。
\end{itemize}
\end{block}
\end{column}
\end{columns}

\sourcenote{snapshot/analysis_summary.md}
\end{frame}

\begin{frame}{Headline 对比表}
\scriptsize
{\setlength{\tabcolsep}{3pt}
\begin{tabular}{L{0.34\linewidth}R{0.09\linewidth}R{0.09\linewidth}R{0.09\linewidth}R{0.09\linewidth}R{0.09\linewidth}R{0.09\linewidth}}
\toprule
Model & RMSE & MAE & Bias & P95 & 1\(\sigma\) cov. & 2\(\sigma\) cov. \\
\midrule
Hiroyasu--Arai calibrated & 9.974 & 7.686 & +0.625 & 20.493 & 0.536 & 0.886 \\
Naber--Siebers delay & 8.840 & 6.780 & -0.344 & 17.922 & 0.620 & 0.893 \\
Stage-3 MLP \(\Delta P^{0.25}\), seed 42 & 6.696 & 4.875 & -1.149 & 14.118 & 0.745 & 0.979 \\
\good{Stage-3 MLP \(\Delta P^{0.5}\), 5-seed mean} & \good{6.211} & \good{4.587} & \good{-0.838} & \good{13.029} & \good{0.774} & \good{0.981} \\
Sparse heteroscedastic GP, 5-seed mean & 6.898 & 5.206 & -1.184 & 14.365 & 0.670 & 0.950 \\
\bottomrule
\end{tabular}
}

\vspace{0.7em}
\small
\begin{block}{写进 thesis 时的定位}
旧 MLP 结果不再作为神经网络 headline，只作为 \(\Delta P\) 指数从
0.25 改到 0.5 之前的 reference row。GP 是强 ML baseline：校准接近 nominal，
但精度不及新 MLP 五 seed 均值。
\end{block}

\sourcenote{snapshot/headline_comparison.csv}
\end{frame}

\begin{frame}{物理 baseline I：Hiroyasu--Arai / Zhou}
\small
\begin{block}{理论角色}
\footnotesize
Hiroyasu--Arai 把喷雾穿深写成分阶段经验律：早期受喷出速度控制，破碎后转入
entrainment-dominated 的平方根时间标度。Zhou 型后喷射扩展进一步给出 EOI 后
\((t-t_i)^{1/4}\) 分支。
\end{block}

\vspace{-0.2em}
\[
S(t)=
\begin{cases}
k_v\sqrt{2\Delta P/\rho_l}\,t, & t<t_b,\\[0.25em]
k_p(\Delta P/\rho_a)^{1/4}d_n^{1/2}t^{1/2}, & t\ge t_b.
\end{cases}
\]
\[
t_b=k_{bt}\frac{\rho_l d_n}{\sqrt{\rho_a\Delta P}},
\qquad
S_{\mathrm{post}}\propto
(\Delta P/\rho_a)^{1/4}d_n^{1/2}t_i^{1/4}(t-t_i)^{1/4}.
\]

\footnotesize
本研究使用 calibrated 版本作为强物理 baseline；其 uncertainty 来自残差统计，
不是输入相关的 learned \(\sigma(x,t)\)。
\end{frame}

\begin{frame}{物理 baseline II：Naber--Siebers}
\small
\begin{block}{理论角色}
\footnotesize
Naber--Siebers 是 mixing-limited penetration 相关式。它把主要工况依赖折叠进
物理尺度 \(A\)，再乘以 \(\sqrt{t}\) 型时间增长。
\end{block}

\[
S(t)=K\,
\underbrace{\left(\frac{\Delta P}{\rho_a}\right)^{1/4}d_n^{1/2}}_{A_{\mathrm{NS}}}
\sqrt{\max(t-t_d,0)}.
\]

\begin{itemize}
  \item \(K\)：global scale；\(t_d\)：液压/观测延迟。
  \item optional cone-angle variant 可使用 video-measured \(\theta\)，但 headline
  对比使用无需 oracle visual feature 的 delay variant。
  \item 方法意义：给 MLP 一个可解释、低自由度、物理刚性的 baseline。
\end{itemize}
\end{frame}

\begin{frame}{ML baseline：Sparse heteroscedastic GP}
\small
\begin{block}{为什么需要 GP}
\footnotesize
Gaussian Process 是小数据、平滑函数和 uncertainty 的强 baseline。若 GP 同时赢
accuracy 与 calibration，MLP 的必要性会被削弱；若 MLP 赢 accuracy 而 GP 赢
calibration，就形成清楚的 trade-off。
\end{block}

\[
f(x)\sim\mathcal{GP}(m(x),k_\theta(x,x')),
\qquad
\hat S = S/A_{\mathrm{scale}}.
\]

\begin{itemize}
  \item 使用 sparse variational GP 避免 exact GP 的 \(O(n^3)\) 成本。
  \item 输入与 MLP 相同：time、tilt、plume count、injection duration、control backpressure。
  \item Heteroscedastic two-stage 版本：mean GP 预测 \(\mu\)，log-variance GP
  从 residual scale 学 \(\sigma(x,t)\)。
  \item 读法：GP 接近 nominal coverage，但 external clean accuracy 不及
  \(\Delta P^{0.5}\) MLP。
\end{itemize}
\end{frame}

\begin{frame}{GP 能打赢 MLP 吗？}
\small
\begin{columns}[T,onlytextwidth]
\begin{column}{0.48\linewidth}
\begin{block}{Accuracy：没有}
{\setlength{\tabcolsep}{3pt}
\begin{tabular}{L{0.48\linewidth}R{0.20\linewidth}R{0.20\linewidth}}
\toprule
Metric & MLP & GP \\
\midrule
RMSE & \good{6.211} & 6.898 \\
MAE & \good{4.587} & 5.206 \\
P95 & \good{13.029} & 14.365 \\
\bottomrule
\end{tabular}
}

\vspace{0.4em}
\footnotesize
相对新 MLP，GP 的 RMSE/MAE/P95 分别高 11.1\%/13.5\%/10.2\%。
\end{block}
\end{column}

\begin{column}{0.48\linewidth}
\begin{block}{Calibration：GP 更接近 nominal}
{\setlength{\tabcolsep}{3pt}
\begin{tabular}{L{0.48\linewidth}R{0.20\linewidth}R{0.20\linewidth}}
\toprule
Metric & MLP & GP \\
\midrule
1\(\sigma\) cov. & 0.774 & \good{0.670} \\
2\(\sigma\) cov. & 0.981 & \good{0.950} \\
\bottomrule
\end{tabular}
}

\vspace{0.4em}
\footnotesize
Gaussian nominal 为 0.683/0.954。GP 的不确定性更贴近 nominal，MLP 在 clean
set 上偏保守。
\end{block}
\end{column}
\end{columns}

\vspace{0.2em}
\begin{block}{结论}
GP 是强 baseline，尤其证明 uncertainty comparison 是公平的；但在同一 external
clean 协议下，它没有打赢 \(\Delta P^{0.5}\) Stage-3 MLP。
\end{block}

\sourcenote{snapshot/headline_comparison.csv}
\end{frame}

\begin{frame}{OOD / LONO：方法含义}
\small
\begin{block}{Leave-one-nozzle-out protocol}
\footnotesize
LONO 将某个 nozzle family / experiment label 的全部样本从训练集中移除，
只在该 held-out group 上测试。它比普通 group-aware split 更严格，因为测试喷嘴
在训练阶段完全不可见。
\end{block}

\begin{block}{它回答的问题}
\footnotesize
External clean eval 主要验证同分布条件下的性能；LONO 验证模型能否在未见喷嘴族上
保持物理 scaling 与 residual learning 的有效性。这里的统计对象是 fold-to-fold
variation，而不是 seed-to-seed variation。
\end{block}

\begin{block}{当前 5-fold 结果}
\footnotesize
MLP \(\Delta P^{0.5}\)：RMSE \(7.435\pm0.692\) mm；HA：
\(10.713\pm1.542\) mm；NS：\(9.784\pm2.102\) mm。相对 HA/NS 的 RMSE
下降约 30.6\% / 24.0\%。
\end{block}
\end{frame}

\begin{frame}{LONO 结果：每个 held-out nozzle 的 RMSE}
\centering
\includegraphics[width=0.92\linewidth,height=0.72\textheight,keepaspectratio]{lono_rmse_by_fold.png}

\vspace{0.15em}
\footnotesize
五个 held-out fold 中，MLP 均低于 HA 与 NS；最难 fold 为 BC20241017\_HZ\_Nozzle2，
MLP RMSE 为 \SI{8.517}{mm}，仍低于 HA \SI{11.754}{mm} 与 NS \SI{11.236}{mm}。

\sourcenote{MLP/runs\_mlp/lono\_pipeline\_20260509\_215613/per\_fold.csv}
\end{frame}

\begin{frame}{LONO 结果：coverage 是否随 OOD 崩掉？}
\centering
\includegraphics[width=0.92\linewidth,height=0.72\textheight,keepaspectratio]{lono_coverage_by_fold.png}

\vspace{0.15em}
\footnotesize
MLP 的 1\(\sigma\) coverage 均值为 0.701，接近 Gaussian nominal 0.683；
Nozzle2 fold 明显偏 under-coverage，说明该喷嘴族仍是 OOD 风险点，而不是被平均值掩盖。

\sourcenote{MLP/runs\_mlp/lono\_pipeline\_20260509\_215613/aggregate.csv}
\end{frame}

\begin{frame}{误差指标：物理 baseline、GP 与 MLP}
\centering
\includegraphics[width=0.90\linewidth,height=0.70\textheight,keepaspectratio]{metric_bars_rmse_mae_p95.png}

\vspace{0.2em}
\footnotesize
在该 clean 评估协议下，sparse GP 明显优于物理 baseline，但仍未超过
\(\Delta P^{0.5}\) Stage-3 MLP 的 RMSE、MAE 和尾部 P95 error。

\sourcenote{snapshot/headline_comparison.csv}
\end{frame}

\begin{frame}{五 seed 外部评估：MLP vs GP}
\centering
\includegraphics[width=0.92\linewidth,height=0.72\textheight,keepaspectratio]{per_seed_rmse.png}

\vspace{0.1em}
\footnotesize
\(\Delta P^{0.5}\) MLP 五 seed RMSE 均值为 \SI{6.211}{mm}；GP 五 seed
external mean 为 \SI{6.898}{mm}。GP 的最好 seed 仍高于 MLP 的 5-seed mean。

\sourcenote{snapshot/new_mlp_5seed_aggregate.csv}
\end{frame}

\begin{frame}{Coverage：GP 更接近 nominal，MLP 更准}
\centering
\includegraphics[width=0.90\linewidth,height=0.72\textheight,keepaspectratio]{coverage_comparison.png}

\vspace{0.1em}
\footnotesize
新 MLP 的 1\(\sigma\)/2\(\sigma\) coverage 为 0.774/0.981，属于 clean set 上
over-coverage；GP 为 0.670/0.950，更接近 Gaussian nominal 0.683/0.954。
因此 GP 赢在 calibration closeness，MLP 赢在 accuracy。

\sourcenote{snapshot/headline_comparison.csv}
\end{frame}

\begin{frame}{写入 thesis 的推荐表述}
\small
\begin{block}{Abstract / introduction 版本}
\(\Delta P^{0.5}\) feature engineering reduced the Stage-3 MLP RMSE from
\SI{6.70}{mm} for the previous \(\Delta P^{0.25}\) seed-42 model to
\SI{6.21}{mm} \(\pm\) \SI{0.28}{mm} over five full-clean external evaluations.
\end{block}

\begin{block}{Results 版本}
Under the same \num{795561}-point clean diagnostic protocol, the
\(\Delta P^{0.5}\) Stage-3 MLP reduced RMSE by 37.7\% relative to the calibrated
Hiroyasu--Arai correlation, by 29.7\% relative to the Naber--Siebers delay
model, and by 10.0\% relative to a sparse heteroscedastic GP.
\end{block}

\sourcenote{snapshot/DATA_SOURCES.md}
\end{frame}

\begin{frame}{GP 与稳定性 caveat}
\small
\begin{block}{GP baseline 版本}
\footnotesize
The sparse heteroscedastic GP reached near-nominal uncertainty coverage
(0.670/0.950 for 1\(\sigma\)/2\(\sigma\)), but its external-clean RMSE, MAE, and
P95 error were 11.1\%, 13.5\%, and 10.2\% higher than the \(\Delta P^{0.5}\)
Stage-3 MLP five-seed mean.
\end{block}

\begin{block}{Limitations 版本}
\footnotesize
The five-seed aggregate is a \(\Delta P^{0.5}\) Stage-3 pipeline aggregate, not a
strict same-ablation seed study: seed 2024 selected \texttt{raw\_reliable\_no\_kd},
whereas the other four seeds selected \texttt{anchor\_off}.
\end{block}

\sourcenote{snapshot/DATA_SOURCES.md}
\end{frame}

\begin{frame}{数据来源与更新入口}
\scriptsize
\begin{tabular}{L{0.24\linewidth}L{0.68\linewidth}}
\toprule
Item & Path \\
\midrule
Frozen snapshot &
\sourcepath{Thesis/generated/baseline_comparison_20260509/} \\
Source report &
\sourcepath{MLP/baseline/comparison_reports/gp_vs_mlp_20260509/} \\
Headline table &
\sourcepath{Thesis/generated/baseline_comparison_20260509/headline_comparison.csv} \\
Five-seed aggregate &
\sourcepath{Thesis/generated/baseline_comparison_20260509/new_mlp_5seed_aggregate.csv} \\
Per-seed eval rows &
\sourcepath{Thesis/generated/baseline_comparison_20260509/per_seed_new_mlp_eval.csv} \\
GP bootstrap &
\sourcepath{Thesis/generated/baseline_comparison_20260509/gp_bootstrap_summary.json} \\
LONO result &
\sourcepath{MLP/runs_mlp/lono_pipeline_20260509_215613/} \\
Figure builder &
\sourcepath{Thesis/generated/baseline_comparison_20260509/make_figures.py} \\
\bottomrule
\end{tabular}

\vspace{0.6em}
\begin{block}{刷新流程}
\footnotesize
更新 comparison report 的 CSV \(\rightarrow\) 覆盖本 snapshot 中的 CSV
\(\rightarrow\) 运行
\sourcepath{.venv/Scripts/python.exe Thesis/generated/baseline_comparison_20260509/make_figures.py}
\(\rightarrow\) 重编译本 slides。
\end{block}

\begin{block}{协议标签}
\footnotesize
所有 headline rows 均为 external clean eval：
\num{795561} points，\num{33845} trajectories。
\end{block}
\end{frame}

\end{document}
